\chapter{实验与分析}
\label{chap:eval}

\section{实验环境}

\begin{table}[htbp]
  \centering
  \caption{实验环境}
  \label{tab:env}
  \begin{tabular}{ll}
    \toprule
    项 & 配置 \\
    \midrule
    基线 & xv6-riscv（课题分支） \\
    模拟器 & QEMU \texttt{virt}，\texttt{CPUS=1} \\
    后端 & EMUL（8\,MiB）/ MAP / PMEM2（见第~\ref{chap:design}~章） \\
    回归 & \texttt{usertests}、\texttt{contract}/\texttt{contract-strict}/MAP/PMEM2（CI） \\
    计时 & 内核 \texttt{r\_time()} 相对单位；每 $N$ 取 5 trial 中位数；每 $N$ 独立开机 \\
    \bottomrule
  \end{tabular}
\end{table}

数据复现：\texttt{make bench-sweep} 回填 \texttt{docs/BENCH.md}；契约：\texttt{make contract}、\texttt{contract-strict}、\texttt{contract-pmem}、\texttt{contract-pmem2}。公平消融见 \texttt{docs/ABLATION.md}。语义与故障矩阵见第~\ref{sec:semantics}~节与 \texttt{docs/SEMANTICS.md}。

\textbf{两套结论：}架构面——相对 xv6 日志路径的量级差；公平面——同载荷消融效应量。禁止把 \texttt{fratio} 读成同等 POSIX API 加速比。

\section{功能、崩溃与掉电（尖刺证据）}

\begin{table}[htbp]
  \centering
  \caption{契约验收用例}
  \label{tab:contract}
  \begin{tabular}{llll}
    \toprule
    用例 & 命令/阶段 & 期望 & 说明 \\
    \midrule
    总闸 & \texttt{make contract} & \texttt{ONBOX\_PMFS\_PASS} & EMUL 必过子集 \\
    严格 & \texttt{contract-strict} & 同上 + A0--A4/\texttt{FAULT\_*} & POS/NEG/效应量 \\
    MAP & \texttt{contract-pmem} & \texttt{mode=MAP}+PMEM1 & 窗口映射 \\
    掉电 & \texttt{contract-pmem2} & \texttt{PMEM2\_PASS} & 文件后端杀启 \\
    remount & \texttt{REMOUNT} & \texttt{equal=1} & sync 后字节等价 \\
    故障 & \texttt{FAULT\_*} & 不可见 / chk 拒绝 & 未 sync；半页撕坏 \\
    POSIX & \texttt{pmfs\_posix} & \texttt{POSIX\_PASS} & 真实文件 API，不升主线 \\
    \bottomrule
  \end{tabular}
\end{table}

分析：契约套件使「完好 SB / 未提交不可见 / remount 等价 / 半页拒绝 / 窗口 persist」可机器判定；相对仅演示双槽结构，这是方法学差异，构成本文尖刺。

\section{性能：随 $N$ 变化的延迟}

对比路径故意不对等（表~\ref{tab:unequal}）。ratio$\gg 1$ 主要反映避免 virtio 日志与块栈的架构差，\textbf{不可}解读为同等完整 POSIX 下的千倍加速。

\begin{table}[htbp]
  \centering
  \caption{故意不对等的对比维度}
  \label{tab:unequal}
  \begin{tabular}{p{2.0cm}p{5.0cm}p{5.2cm}}
    \toprule
    维度 & PMFS-KV / File & xv6 fs+log \\
    \midrule
    API & KV：put$\times N$+seal+get；
         File：create/write/fsync$\times N$
         & open/create+write(8B)$\times N$+sync \\
    持久语义 & 双槽 SB + inode/dentry 页 & 块日志检查点 \\
    I/O 栈 & 字节寻址窗口 / EMUL & virtio-blk + bio + log \\
    元数据 & 有限 inode/dentry（无块位图日志） & 目录+inode+位图+log \\
    \bottomrule
  \end{tabular}
\end{table}

\begin{table}[htbp]
  \centering
  \caption{EMUL 下随 $N$ 变化的中位数与相对吞吐（5 trial，CDC 融合后）}
  \label{tab:bench-n}
  \begin{tabular}{rrrrrrr}
    \toprule
    $N$ & med\_kv & med\_file & med\_xv6 & fratio$\times$100 & thr\_file & wa\_x100 \\
    \midrule
    8  & 33507 & 304451 & 5684939  & 1867 & 26 & 15384 \\
    16 & 38448 & 486211 & 11208338 & 2305 & 32 & 17889 \\
    32 & 53065 & 659056 & 25648605 & 3891 & 48 & 17721 \\
    \bottomrule
  \end{tabular}
\end{table}

\FigNLatency

观察：（1）xv6 随 $N$ 近似线性；（2）PMFS-File 每文件 fsync 仍避开块日志栈，fratio$\gg 1$；（3）File 路径因 CDC/stream/meta 提交高于纯 KV，但仍与 xv6 差一个数量级以上；（4）结论写作\textbf{量级优势 + 边界}。

\section{写放大}

\FigWaN

小 value（8\,B）相对 4\,KiB 页对齐主导放大。CDC 融合后本负载下 $\mathrm{wa\_x100}$ 维持在高位且随 $N$ 近似平台（表~\ref{tab:bench-n}），属页粒度与 run 打包的预期行为；与 xv6 日志放大不同源，不可横向比绝对值。

\section{MAP 与 EMUL 对照}

\begin{table}[htbp]
  \centering
  \caption{$N{=}32$ 时 EMUL 与 MAP}
  \label{tab:map-emul}
  \begin{tabular}{lrrrr}
    \toprule
    mode & med\_kv & med\_file & med\_xv6 & fratio$\times$100 \\
    \midrule
    EMUL & 53065  & 659056  & 25648605 & 3891 \\
    MAP  & 113804 & 4301198 & 24439543 & 568 \\
    \bottomrule
  \end{tabular}
\end{table}

\FigMapEmul

\FigThr

文件路径相对吞吐在 EMUL 下稳定高于 xv6 一个数量级以上（表中 \texttt{thr\_file}）。MAP 抬升软件路径开销后 fratio 下降但仍 $\gg 1$。EMUL 仍\textbf{不}代表真 NVDIMM 延迟。

\section{公平消融（公平面）}

\begin{table}[htbp]
  \centering
  \caption{A0--A4 公平消融效应量（EMUL，\texttt{contract-strict} 实测）}
  \label{tab:ablation}
  \begin{tabular}{lccrrr}
    \toprule
    变体 & CDC & batch seal & cycles & wa\_x100 & 解读 \\
    \midrule
    A0 & 开 & 开 & 30029 & 1600 & 默认策略 \\
    A1 & 关 & 开 & 27490 & 1600 & 定长+批量 \\
    A2 & 开 & 关 & 52066 & 4800 & 每 chunk seal，WA/时延升 \\
    A3 & 关 & 关 & 36238 & 3200 & 定长+立即 seal \\
    A4 & 开 & 开 & 24146 & 1600 & A0 复测（稳定性） \\
    \bottomrule
  \end{tabular}
\end{table}

同 payload 下，关闭批量 seal（A2/A3）显著抬升 $\mathrm{wa\_x100}$ 与 cycles；A0/A4 同配置复测同量级，说明效应量可重复。公平面只据此表陈述，\textbf{不}与表~\ref{tab:bench-n} 的 fratio 混读。

\section{实现难点与失败史}

下列踩坑用于证明系统是「调出来的」，并可复核。

\paragraph{P0 启动栈溢出}
\texttt{main()} 跑在原 4\,KiB/hart 的 \texttt{stack0} 上；自测栈上放置多页缓冲导致溢出，污染邻近内存。修复：页缓冲改 \texttt{static}；\texttt{stack0} 扩至 16\,KiB/hart。

\paragraph{P1 \texttt{rdcycle} 非法指令}
\texttt{timerinit} 未开放 \texttt{CY} 时，S 态 \texttt{rdcycle} 触发 \texttt{kerneltrap}。修复：改用已允许的 \texttt{r\_time()}，并在论文中不作墙钟 $\mu$s 解读。

\paragraph{P2 QEMU NUMA / RAM 尺寸不匹配}
仅配置 160\,M 文件后端而默认 $-m$ 仍为 128\,M 时，QEMU 立即退出，契约表现为超时。修复：文件后端启动同步 $-m$\,160M。

\paragraph{P3 MAP 在 $-m$\,128M 下写 PMEMBASE}
无物理后备窗口上的写探测导致 store fault。修复：仅 \texttt{-DPMEM\_MAP} 时映射并探测。

\paragraph{P4 Bench 过早杀进程截断契约行}
匹配到 \texttt{BENCH\_PASS} 子串即杀 QEMU，导致字段截断。修复：等待 \texttt{median\_xv6=} / \texttt{phys=} / \texttt{wa\_x100=} 齐全。

\paragraph{P5 多 trial 撑爆段数}
连续 seal 耗尽段视图导致后续 trial 失败。修复：提高 run/段余量、attach 压力阀、\textbf{每 $N$ 独立开机}扫参。

\paragraph{P6 题面「文件系统」与实现漂移}
早期曾偏并行 KV。修复：升格 dir/extent/CDC + 深化 \texttt{pmfs\_posix}，并以语义清单写死差集，避免题实不符。

\section{局限与威胁效度}

\begin{enumerate}[leftmargin=2em]
  \item QEMU 非真 NVDIMM；附录给出扩展路径但不喧宾夺主。
  \item 架构面性能对比 API 异构；结论限于量级说明。公平面仅 A0--A4 同载荷可比。
  \item 文件语义为可测子集：无 \texttt{lseek}/\texttt{mmap}/\texttt{rename}；容量受 extent/chunk 上限约束；默认 \texttt{CPUS=1}。
  \item 计时为相对单位；外推墙钟或真机延迟需重新校准。
\end{enumerate}

\section{本章小结}

尖刺（契约套件，含 remount/故障）+ 支撑柱（文件子系统/KV 核/PMEM2）+ 两套结论口径 + 可复核踩坑，构成实验章叙事主线；性能图服务于量级与写放大机制说明，消融表服务于公平面效应量，而非绝对公平微基准夺冠。
